\documentclass[11pt, a4paper]{article}

\usepackage[a4paper, top=2.5cm, bottom=2.5cm, left=2cm, right=2cm]{geometry}
\usepackage{fontspec}

\usepackage[english]{babel}
\usepackage{graphicx}
\usepackage{float}
\usepackage{caption}
\usepackage{amsmath, amssymb, amsthm}
\usepackage{booktabs}
\usepackage{enumitem}
\usepackage[numbers]{natbib}
\usepackage[colorlinks=true, linkcolor=blue, citecolor=red, urlcolor=blue]{hyperref}
\usepackage{authblk}

\setmainfont{Noto Serif}
\captionsetup{font=small, labelfont=bf}

% --- THEOREMS ---
\theoremstyle{plain}
\newtheorem{theorem}{Theorem}[section]
\newtheorem{lemma}[theorem]{Lemma}
\newtheorem{proposition}[theorem]{Proposition}
\newtheorem{corollary}[theorem]{Corollary}

\theoremstyle{definition}
\newtheorem{definition}[theorem]{Definition}

\theoremstyle{remark}
\newtheorem{remark}[theorem]{Remark}

% --- TITLE ---
\title{\textbf{Liouville-Type Results for Infinity Elliptic Equations with Gradient-Dependent and Hardy--Hénon Nonlinearities}}

\author[1,2]{Tran Minh Hau}
\author[2,3]{Huynh Trung Hieu}

\affil[1]{University of Information Technology, VNU-HCM, Vietnam}
\affil[2]{Vietnam Institute for Advanced Study in Mathematics (VIASM)}
\affil[3]{Ho Chi Minh City University of Education}

\date{\small Advisors: Prof. Hai-Long Dao, Prof. Dac-Tuan Ngo, Dr. Tuan-Minh Nguyen \\ August 2025}

\begin{document}

\maketitle

\begin{abstract}
This paper establishes Liouville-type nonexistence theorems for infinity elliptic equations involving gradient-dependent nonlinearities and Hardy--Hénon weights. 

Within the viscosity solution framework, we identify a dimension-free critical exponent $p = 3$ arising from the cubic structure of the infinity Laplacian, and prove the collapse of nontrivial solutions in the subcritical regime.

We further show that a sharp gradient threshold $q = 3$ governs the transition to gradient-dominated behavior, leading to complete solution collapse under mild decay assumptions.

These results contribute to the qualitative theory of highly degenerate elliptic equations and provide a theoretical bridge to gradient regularization in modern machine learning.
\end{abstract}

% =========================
\section{Introduction}

Liouville-type theorems play a central role in nonlinear elliptic theory \cite{gidas1981}. 
The classical Lane--Emden equation
\[
-\Delta u = u^p
\]
has well-known critical exponents depending on dimension.

Extensions such as the Hardy--Hénon equation introduce spatial weights \cite{le2025}. 

The infinity Laplacian
\[
\Delta_\infty u = \langle D^2 u \nabla u, \nabla u \rangle
\]
is highly degenerate and requires viscosity solutions \cite{crandall1992, aronsson1967}.

We study:
\[
-\Delta_\infty u + c|\nabla u|^q = \lambda |x|^a u^p
\]

% =========================
\section{Main Results}

\begin{theorem}
Let $u$ be a non-negative viscosity solution of
\[
-\Delta_\infty u = \lambda |x|^a u^p.
\]
If $1 < p \le 3$, then $u \equiv 0$.
\end{theorem}

\begin{theorem}
Let $u$ satisfy
\[
-\Delta_\infty u + c|\nabla u|^q \ge \lambda |x|^a u^p.
\]
If $q > 3$ and $p < q$, then $u \equiv 0$.
\end{theorem}

% =========================
\section{Proof Strategy}

The proofs rely on:
\begin{itemize}
\item scaling invariance of the infinity Laplacian
\item radial reduction techniques
\item viscosity solution comparison principles
\end{itemize}

See \cite{crandall1992, crandall2001, peres2009}.

% =========================
\section{Connections to Machine Learning}

The infinity Laplacian corresponds to minimizing the Lipschitz constant \cite{crandall2001}. 

Gradient regularization in deep learning \cite{finlay2018, gao2022} can be interpreted as analogous to PDE constraints:
\[
\mathcal{L}(f) = \mathbb{E}[\text{loss}] + c \mathbb{E}[|\nabla f|^q]
\]

Our results suggest that excessive gradient penalization leads to collapse of solutions, analogous to over-regularization in neural networks.

% =========================
\section{Conclusion}

We established Liouville-type nonexistence results and identified sharp thresholds $p=3$ and $q=3$ governing solution behavior.

Future work includes parabolic extensions and anisotropic weights.

% =========================
\bibliographystyle{unsrt}
\bibliography{references}

\end{document}
